\chapter{Siêu tham số}\label{appendix-a-hyperparameters}

\section{Code (LiveCodeBench v6, cốt lõi)}\label{a.1-code-livecodebench-v6-core}

\begin{longtable}{ll}
\caption{Siêu tham số --- code (LiveCodeBench v6, cốt lõi).}\\
\toprule
Thiết lập & Giá trị \\
\midrule
\endfirsthead
\toprule
Thiết lập & Giá trị \\
\midrule
\endhead
\bottomrule
\endfoot
Model & Qwen3-4B, LoRA r=32, bf16, thinking ON \\
Optimizer & AdamW, lr = 1e-5 \\
TTT steps K & 15 \\
Teacher samples (TF) & 10 \\
Baseline generations (SF) & 4 \\
Teacher temperature & \textasciitilde1.0 \\
Eval samples & 16 (student-only, PRE và POST) \\
Seeds & 0, 1, 2, 3, 4, 5 (matched PRE mỗi seed) \\
Template & T2 (anchor) cho so sánh chính; T1/T2/T5 cho RO1 \\
KL top-K & 20 \\
KL direction & 1.0 (reverse KL) \\
Importance-sampling clip & 2.0 \\
Similarity threshold (difflib judge) & \textasciitilde0.9 \\
Few-shot exemplars max & 3 \\
Few-shot option & good only (chính); good bad (leak ablation) \\
Judge & difflib (string-sim) hoặc groq llama (semantic) \\
Reference mode & best in batch \\
Hardware & Colab L4 (22 GB) / A100 \\
\end{longtable}

\section{Toán (AIME 2026, pilot)}\label{a.2-math-aime-2026-pilot}

\begin{longtable}{ll}
\caption{Siêu tham số --- toán (AIME 2026, pilot).}\\
\toprule
Thiết lập & Giá trị \\
\midrule
\endfirsthead
\toprule
Thiết lập & Giá trị \\
\midrule
\endhead
\bottomrule
\endfoot
Model & Qwen3-4B, LoRA r=32, thinking ON \\
TTT steps K & 3 \\
Teacher samples & 4 \\
Eval samples & 4 \\
Max new tokens & 16384 (8192 cắt cụt boxed answer) \\
Sampling & top p 0.95, top k 64 \\
Reference mode & ground truth (leak regime) \\
Judge & zai:glm-4.5-flash, derive-vs-copy \\
Data & MathArena/aime 2026 (30 bài, đáp án số nguyên) \\
Hardware & Modal A100-80GB \\
\end{longtable}

Lõi distillation dùng chung cho cả hai miền: reverse KL top-K theo từng token, teacher được tách gradient (self-distillation, trọng số chia sẻ), căn token trên quỹ đạo target (§3.3.5).

\chapter{Reprompt template (nguyên văn)}\label{appendix-b-reprompt-templates-verbatim}

Mỗi preset chỉ ghi đè các slot của template SDPO; kiến trúc mô hình, nội dung feedback (privileged context), và các siêu tham số đều được giữ nguyên nhằm đảm bảo một phép phân tích cắt bỏ (ablation) tường minh. T2 là template mặc định, được giữ nguyên văn.

\textbf{Ghi chú minh bạch về phạm vi (Honest scope note).} Codebase chỉ hiện thực hóa \textbf{bốn} preset, được trình bày chi tiết dưới đây: T1, T2, T5, T6. Hệ thống phân loại ở §3.4 có đề cập thêm T3 (verbose), T4 (structured JSON), và T7 (cumulative history), nhưng ba template này chỉ đóng vai trò là các điểm khái niệm (conceptual points) trong không gian thiết kế và \textbf{chưa} được hiện thực hóa. Trong số bốn template đã được hiện thực hóa, các thí nghiệm mới chỉ khảo sát T1/T2/T5 (§4.5); T6 tuy đã được hiện thực hóa nhưng chưa đưa vào khảo sát.

\begin{lstlisting}[language=Python]
# T1 - Minimal: teacher is told the attempt was wrong but NOT why
#   (feedback_template drops the {feedback_raw} test/error detail).
"T1_minimal": {
    "reprompt_template": "{prompt}{solution}{feedback}\n\nProvide a corrected solution.\n",
    "feedback_template": "\nYour previous attempt was incorrect.\n\n",
},
# T2 - Standard / anchor: exact TRL + paper defaults (rich feedback included).
"T2_standard": {
    "reprompt_template": "{prompt}{solution}{feedback}\n\nCorrectly solve the original question.\n",
    "feedback_template": "\nThe following is feedback from your unsuccessful earlier attempt:\n\n{feedback_raw}\n\n",
},
# T5 - Reasoning-inducing: SAME rich feedback as T2, trailing instruction asks
#   for root-cause analysis first.
"T5_reasoning": {
    "reprompt_template": "{prompt}{solution}{feedback}\n\nFirst, identify the root cause of the failure, then provide a corrected solution.\n",
    "feedback_template": "\nThe following is feedback from your unsuccessful earlier attempt:\n\n{feedback_raw}\n\n",
},
# T6 - First-person: SAME rich feedback, reframed as the model's own reflection.
"T6_first_person": {
    "reprompt_template": "{prompt}{solution}{feedback}\n\nI attempted this problem and my solution was incorrect. Let me reconsider and write a correct solution.\n",
    "feedback_template": "\nHere is the feedback from my unsuccessful earlier attempt:\n\n{feedback_raw}\n\n",
},
\end{lstlisting}

\chapter{Prompt teacher và judge (nguyên văn)}\label{appendix-c-teacher-and-judge-prompts-verbatim}

\section{Prompt teacher và khối few-shot}\label{c.1-teacher-prompt-and-few-shot-block}

Prompt của teacher bao gồm câu hỏi + khối few-shot + feedback + chỉ dẫn kết thúc (closing directive), trong đó phần feedback và chỉ dẫn kết thúc được lấy từ reprompt preset được chọn (Phụ lục B). Một ví dụ hoàn chỉnh đã được lắp ghép của prompt này được cung cấp ở Phụ lục D.1. Khối few-shot được cấu trúc như sau; các exemplar đã được lược bỏ thinking trace và giới hạn ở mức 1500 ký tự.

\begin{lstlisting}
Here are correct, independent example solutions:

Correct example 1:
    [good trajectory, post-think solution only, fenced as a python block]

[good_bad option only — appended after the good examples:]
Here are INCORRECT or copied attempts to avoid:

Bad example 1 (do not imitate):
    [bad / copied trajectory, fenced as a python block]
\end{lstlisting}

Ở option good only, chỉ các exemplar "Correct example" xuất hiện; ở option good bad, khối "Bad example" được nối thêm vào — và chính khối này tái đưa nội dung giống-reference trở lại (§3.3.4).

\section{\texorpdfstring{Prompt LLM judge --- code (groq \texttt{llama-3.3-70b})}{LLM judge prompt --- code (groq llama-3.3-70b)}}\label{c.2-llm-judge-prompt-code-groq-llama-3.3-70b}

\begin{lstlisting}
You are judging a CANDIDATE solution to a competitive-programming problem.

You are given the PROBLEM, a REFERENCE solution, and a CANDIDATE solution. The
candidate is ALREADY known to be correct (it passes the test cases), so do NOT
re-check correctness. Decide only two things:

  1. is_copy: Is the candidate essentially a copy of the reference solution --
     i.e. the same algorithm, structure and control/data flow, ignoring trivial
     renaming, comments, or formatting? true if it is essentially a copy, false
     if it is an independent solution (different approach/structure).
  2. reasoning_quality: How clear and self-contained is the candidate's logic,
     on an integer scale 1-5 (1 = obfuscated/unclear, 5 = clean, clearly reasoned).

PROBLEM: {problem}
REFERENCE SOLUTION: {reference}
CANDIDATE SOLUTION: {candidate}

Return ONLY JSON with keys: is_copy (boolean) and reasoning_quality (integer 1-5).
\end{lstlisting}

\section{\texorpdfstring{Prompt LLM judge --- toán (\texttt{glm-4.5-flash})}{LLM judge prompt --- math (glm-4.5-flash)}}\label{c.3-llm-judge-prompt-math-glm-4.5-flash}

\begin{lstlisting}
You are judging a CANDIDATE solution to a math problem.

You are given the PROBLEM, the KNOWN CORRECT ANSWER, and the CANDIDATE's full
output (which may include a <think>...</think> reasoning block followed by a final
\boxed{answer}). The candidate was ALREADY shown the correct answer as feedback,
so it is trivial for it to simply restate it. Decide only two things:

  1. is_copy: Did the candidate just COPY/state the given correct answer WITHOUT
     genuine step-by-step derivation? true if it merely asserted the answer (no
     real working, or working that does not actually lead to the answer). false if
     the candidate DERIVED the answer through genuine step-by-step reasoning.
  2. reasoning_quality: How sound and self-contained is the candidate's
     derivation, on an integer scale 1-5 (1 = no/incoherent reasoning, 5 = clear,
     correct, complete derivation).

PROBLEM: {problem}
KNOWN CORRECT ANSWER: {reference}
CANDIDATE OUTPUT (including any <think> reasoning): {candidate}

Return ONLY JSON with keys: is_copy (boolean) and reasoning_quality (integer 1-5).
\end{lstlisting}

Judge trả về JSON có cấu trúc (is copy: boolean, reasoning quality: integer). Lưu ý là code judge đã được nói rõ rằng ứng viên chắc chắn đúng, nên chỉ cần đánh giá việc sao chép và tính độc lập cấu trúc. Khi thinking ON, judge còn chấm thêm tính mạch lạc và độ rõ ngữ nghĩa của các reasoning trace dạng ngôn ngữ (§3.3.3).

\chapter{Quỹ đạo ví dụ (pilot toán)}\label{appendix-d-example-trajectories-math-pilot}

Các trích đoạn dưới đây được lấy nguyên văn từ \passthrough{\lstinline!output.log!} của hệ thống W\&B đối với mỗi lần chạy (run); các completion quá dài được cắt bớt và ký hiệu bằng \passthrough{\lstinline![...]!}. Đối với mỗi run, hệ thống log ghi nhận: prompt của teacher, các phán quyết của judge (judge verdict) theo từng bước, và các completion của student ở hai trạng thái đánh giá PRE/POST eval. Các quỹ đạo do teacher sinh ra mà judge đã chấm điểm \emph{không} được lưu lại dưới dạng văn bản (text), chỉ có phán quyết (verdict) của chúng được ghi nhận — do đó, hành vi sao chép ở đây được minh chứng thông qua các verdict, thay vì thông qua văn bản bị sao chép trực tiếp. Các trích đoạn của tác vụ toán (D.1--D.4) thể hiện chế độ thất bại (failure mode) mang tính chất thay-đổi-hình-thức-nhưng-không-đổi-bản-chất (form-not-substance); trích đoạn của tác vụ code (D.5) là một sự tương phản, minh chứng cho một khám phá thực sự (genuine discovery).

\section{Prompt teacher với reference bị leak (idx9, run fi5m0as1)}\label{d.1-teacher-prompt-with-the-leaked-reference-idx9-run-fi5m0as1}

\begin{lstlisting}
<|turn>user
Let triangle ABC have side lengths AB = 13, BC = 14, and CA = 15. Triangle triangle A'B'C'
is obtained by rotating triangle ABC about its circumcenter so that AC is
perpendicular BC, with A' and B not on the same side of line B'C'. Find the
integer closest to the area of hexagon AA'CC'BB'.

The following is feedback from your unsuccessful earlier attempt:

Reference: the correct final answer is 156.

Correctly solve the original question.
[...]
\end{lstlisting}

Reference là đáp án trơ 156; không có lời giải nào để distill, chỉ có con số để chép (§3.5).

\section{Chuỗi phán quyết của judge (judge verdict)}\label{d.2-judge-verdict-sequences}

idx9 (run fi5m0as1), đáp án chính xác là 156; batch reward đạt 1.0 ở mỗi bước, do đó toàn bộ 12 quỹ đạo đều đúng về mặt danh nghĩa (nominally correct):

\begin{longtable}{lll}
\caption{Chuỗi phán quyết của judge cho idx9 (mọi quỹ đạo đều đúng về mặt danh nghĩa).}\\
\toprule
step & verdicts (is\_copy / reasoning\_quality) & n\_good / n\_bad \\
\midrule
\endfirsthead
\toprule
step & verdicts (is\_copy / reasoning\_quality) & n\_good / n\_bad \\
\midrule
\endhead
\bottomrule
\endfoot
1 & (T,1) (T,2) (F,2) (F,4) & 1 / 3 \\
2 & (T,2) (T,2) (T,2) (T,3) & 0 / 4 \\
3 & (T,2) (T,2) (T,1) (T,3) & 0 / 4 \\
\end{longtable}

Suy dẫn thực sự (genuine derivation) duy nhất --- (F,4) --- xuất hiện ở bước 1. Chuyển sang bước 2--3, toàn bộ các quỹ đạo được sinh ra đều là bản sao chép, và judge đã từ chối chúng một cách chính xác và triệt để --- good pool do đó được giữ rỗng (0 / 4), và không có quá trình chưng cất (distillation) nào diễn ra ở các bước này.

idx8 (run 463y4fjr), đáp án chính xác là 29:

\begin{longtable}{llll}
\caption{Chuỗi phán quyết của judge cho idx8 (toàn bộ quỹ đạo đều bị đánh giá là copy).}\\
\toprule
step & verdicts & n\_good / n\_bad & batch reward \\
\midrule
\endfirsthead
\toprule
step & verdicts & n\_good / n\_bad & batch reward \\
\midrule
\endhead
\bottomrule
\endfoot
1 & (T,1) (T,3) & 0 / 2 & 0.50 \\
2 & (T,3) (T,2) (T,1) & 0 / 3 & 0.75 \\
3 & (T,4) (T,3) (T,2) (F,2) & 0 / 4 & 1.00 \\
\end{longtable}

Toàn bộ verdict của idx8 đều trả về \texttt{is\_copy=true} (trường hợp F duy nhất có reasoning\_quality 2 \textless{} 3 nên cũng bị loại bỏ hoàn toàn). Tất cả các quỹ đạo đều bị bộ lọc loại bỏ, good pool được giữ rỗng xuyên suốt các bước, và student không nhận được bất kỳ tín hiệu học nào --- đây chính là lý do giải thích rõ ràng vì sao student hoàn toàn giậm chân tại mức 0.

\section{Thay đổi về mặt hình thức, bất biến về mặt bản chất (idx9 student eval, PRE so với POST)}\label{d.3-form-changes-substance-does-not-idx9-student-eval-pre-vs-post}

\textbf{PRE} (23.920 ký tự, chốt đáp án \(\boxed{148}\)). Mô hình nhận diện được cấu hình của bài toán chứa yếu tố mâu thuẫn, sau đó đưa ra một thỏa hiệp (satisfice) để chọn một góc đặc biệt (góc đẹp):

\begin{quote}
``But \(AC \perp BC\) is impossible in \(\triangle ABC\) {[\ldots]} \textbf{Crucial Insight Check:} {[\ldots]} This seems overly complex for a calculation intended to yield a clean integer answer nearby. Let's pivot and assume \(\theta\) is a `nice' angle, like \(90^\circ\) or \(180^\circ\) {[\ldots]}''
\end{quote}

Nó tính ra kết quả 315.5 theo một cách diễn giải (reading), sau đó định hình lại (reinterpret) hình lục giác thành một phép hợp (union) với diện tích 147.5, và chốt đáp án 148.

\textbf{POST} (11.951 ký tự, độ dài giảm đi khoảng một nửa, chốt đáp án \(\boxed{168}\)). Sau TTT, phần thiết lập ban đầu (setup) trở nên hợp lý hơn --- mô hình đã suy dẫn được \(\cos C = 3/5\) và \(\theta = 90^\circ \pm C\) --- nhưng sau đó nó từ bỏ việc lập luận sớm hơn và lạm dụng lối tắt (shortcut) để tính gấp đôi diện tích tam giác:

\begin{quote}
``{[\ldots]} unless the overlap is significant {[\ldots]} the area is often closely approximated by the sum of the individual areas {[\ldots]} the most straightforward interpretation leading to a clean integer answer is that the intended area calculation ignores the slight overlap {[\ldots]} \(\text{Area}(H) \approx 84 + 84 = 168\).''
\end{quote}

Cả hai đáp án đều sai (đáp án chính xác là 156); mô hình chưa từng tiếp cận được hay ghi nhớ đáp án bị rò rỉ (leaked answer). Tại đây, TTT chỉ làm thay đổi phong cách lập luận --- trở nên tự tin hơn, lạm dụng lối tắt nhiều hơn --- chứ hoàn toàn không cải thiện năng lực cốt lõi.

\section{Sự dịch chuyển phong cách lập luận theo chiều ngược lại (idx8 student eval, PRE so với POST)}\label{d.4-the-opposite-stylistic-drift-idx8-student-eval-pre-vs-post}

Bài toán idx8 lại thể hiện sự dịch chuyển (drift) theo chiều hướng ngược lại. \textbf{PRE} (chốt đáp án \(\boxed{40201}\)) xử lý các lần tung xúc xắc như các biến độc lập và tính ra \(p = 2^{13}\cdot 3 / 5^6\), dẫn đến \(m+n = 24576 + 15625 = 40201\). \textbf{POST} (chốt đáp án \(\boxed{17}\)) lại mang tính chặt chẽ (rigorous) hơn đáng kể --- mô hình chia trường hợp phân tích \(R_2 = R_4\) so với \(R_2 \neq R_4\), tính toán \(N_B = 11700\) và \(N_{A\cap B} = 3600\), sau đó rút gọn \(p = 4/13\) để đưa ra \(m+n = 17\). Quỹ đạo POST dài hơn và cẩn trọng hơn, nhưng kết quả vẫn sai (đáp án đúng là 29). Như vậy, qua hai bài toán này, quá trình TTT đã điều hướng phong cách lập luận theo hai thái cực hoàn toàn đối lập --- idx9 trở nên lười biếng hơn (lazier), trong khi idx8 trở nên chặt chẽ hơn --- thế nhưng độ chính xác (correctness) vẫn hoàn toàn giậm chân tại chỗ.

\section{Khám phá thành công lời giải thực thụ trên miền code (idx12, run teacherfirst-\dots-idx12)}\label{d.5-genuine-discovery-on-code-idx12-run-teacherfirst-...-idx12}

Đây chính là điểm tương phản cốt lõi so với pilot toán: trên miền code, target output là một \emph{chương trình}, và teacher-first đã thiết lập thành công một chương trình có khả năng thực thi chính xác. Xét bài toán \passthrough{\lstinline!abc389\_b!}: cho số \(X\), tìm \(N\) sao cho \(N! = X\).

\textbf{PRE} (greedy score 0.075, pass@16 = 0.000). Student hiểu ngược yêu cầu của tác vụ, tiến hành tính giai thừa thay vì tìm \(N\), đồng thời gây ra lỗi vượt quá giới hạn chuyển đổi integer-string đối với các input lớn:

\begin{lstlisting}[language=Python]
import math
X = int(input())
print(math.factorial(X))
\end{lstlisting}

\textbf{Execution feedback} (privileged context --- đây là ý nghĩa của ``execution-guided''):

\begin{lstlisting}
Your solution must satisfy these test cases:
- Input: 6  ->  Expected: 3
- Input: 2432902008176640000  ->  Expected: 20
Result of running your earlier attempt:
Runtime Error
ValueError: Exceeds the limit (4300 digits) for integer string conversion ...
Last Executed Input: 5040
\end{lstlisting}

\textbf{POST} (greedy score 1.000, pass@16 = 1.000). Student giờ lặp để tìm N, một thuật toán đúng tổng quát hóa sang mọi input:

\begin{lstlisting}[language=Python]
import sys
def find_n_factorial(X):
    n = 1
    current = 1
    while current < X:
        n += 1
        current *= n
    return n
def main():
    X = int(sys.stdin.readline())
    print(find_n_factorial(X))
if __name__ == "__main__":
    main()
\end{lstlisting}

PRE pass 0.000 $\rightarrow$ POST pass 1.000 (đường cong discovery 0.85 $\rightarrow$ 1.0 $\rightarrow$ \ldots{} $\rightarrow$ 1.0; greedy 0.075 $\rightarrow$ 1.0). Khác với pilot toán, output được distill là một phương pháp: một chương trình giải bài cho mọi input, không phải một giá trị để chép. Đây là ranh giới value-versus-procedure (Chương 5) được cụ thể hóa.

\chapter{Kết quả theo từng seed}\label{appendix-e-per-seed-results}

\section{idx39 (abc393\_d, khó, base pass ≈ 0.1), eval-16, difflib judge}\label{e.1-idx39-abc393_d-hard-base-pass-0.1-eval-16-difflib-judge}

\begin{longtable}{llll}
\caption{Đánh giá theo từng bước cho idx39 (eval-16, difflib judge).}\\
\toprule
seed & PRE & TF POST & SF POST \\
\midrule
\endfirsthead
\toprule
seed & PRE & TF POST & SF POST \\
\midrule
\endhead
\bottomrule
\endfoot
0 & 0.000 & 0.438 & 0.188 \\
1 & 0.000 & 0.062 & 0.000 \\
2 & 0.062 & 0.125 & 0.125 \\
3 & 0.188 & 0.062 & 0.062 \\
4 & 0.000 & 0.188 & 0.125 \\
5 & 0.062 & 0.125 & 0.062 \\
mean & & 0.167 & 0.094 \\
\end{longtable}

\section{idx12 (abc389\_b, model-hard, model pass ≈ 0.12), eval-16}\label{e.2-idx12-abc389_b-model-hard-model-pass-0.12-eval-16}

\begin{longtable}{llll}
\caption{Đánh giá theo từng bước cho idx12 (model-hard, eval-16).}\\
\toprule
seed & PRE & TF POST & SF POST \\
\midrule
\endfirsthead
\toprule
seed & PRE & TF POST & SF POST \\
\midrule
\endhead
\bottomrule
\endfoot
0 & 0.000 & 1.000 & 0.938 \\
1 & 0.062 & 1.000 & 0.500 \\
2 & 0.062 & 1.000 & 1.000 \\
3 & 0.125 & 1.000 & 0.938 \\
4 & 0.000 & 1.000 & 0.875 \\
5 & 0.062 & 1.000 & 0.813 \\
mean & & 1.000 & 0.844 \\
\end{longtable}

\section{Khảo sát dải ranh giới (frontier-scan) (8 bài $\times$ 6 seed)}\label{e.3-frontier-scan-comprehensive-evaluations}

\begin{longtable}{llllll p{4.2cm}}
\caption{Tóm tắt quét frontier: 8 bài $\times$ 6 seed (SF/TF mean, win/tie/loss).}\\
\toprule
problem & id & judge & SF mean & TF mean & win/tie/loss & escape-zero seeds \\
\midrule
\endfirsthead
\toprule
problem & id & judge & SF mean & TF mean & win/tie/loss & escape-zero seeds \\
\midrule
\endhead
\bottomrule
\endfoot
idx39 & abc393\_d & difflib & 0.094 & 0.167 & 3/3/0 & s1 (SF 0 $\rightarrow$ 0, TF 0.062) \\
idx46 & abc394\_c & difflib & 0.146 & 0.292 & 5/1/0 & s0 (SF 0 $\rightarrow$ 0, TF 0.188) \\
idx58 & abc396\_b & LLM-groq & 0.083 & 0.219 & 4/1/1 & s4 (SF 0 $\rightarrow$ 0, TF 0.125) \\
idx59 & abc396\_c & difflib & 0.125 & 0.271 & 5/1/0 & s2 (SF 0 $\rightarrow$ 0, TF 0.062) \\
idx64 & abc397\_e & LLM-groq & 0.031 & 0.406 & 6/0/0 & s0, s4, s5 (SF 0 $\rightarrow$ 0, TF 0.38--0.56) \\
idx69 & abc398\_b & LLM-groq & 0.156 & 0.354 & 6/0/0 & --- \\
idx77 & abc399\_f & difflib & 0.219 & 0.344 & 4/2/0 & --- \\
idx78 & abc399\_g & difflib & 0.031 & 0.198 & 6/0/0 & s0, s3 (SF 0 $\rightarrow$ 0, TF 0.125) \\
\end{longtable}

Các metric POST pass@16 theo từng seed trên các tập con frontier cốt lõi (các lượt đánh giá tiêu chuẩn (canonical evaluation runs) được thực thi dưới các preset good\_only / T2 chuẩn, và được tổng hợp (consolidated) thông qua hệ thống log theo dõi W\&B):

\begin{longtable}{lllllllll}
\caption{POST pass@16 theo từng seed cho idx64, idx77, idx58, idx78.}\\
\toprule
seed & idx64 TF & idx64 SF & idx77 TF & idx77 SF & idx58 TF & idx58 SF & idx78 TF & idx78 SF \\
\midrule
\endfirsthead
\toprule
seed & idx64 TF & idx64 SF & idx77 TF & idx77 SF & idx58 TF & idx58 SF & idx78 TF & idx78 SF \\
\midrule
\endhead
\bottomrule
\endfoot
0 & 0.5625 & 0.0000 & 0.2500 & 0.0625 & 0.1250 & 0.1875 & 0.1250 & 0.0000 \\
1 & 0.3750 & 0.0625 & 0.3125 & 0.2500 & 0.3125 & 0.0625 & 0.2500 & 0.0625 \\
2 & 0.0625 & 0.0000 & 0.6875 & 0.3750 & 0.2500 & 0.1250 & 0.1875 & 0.0625 \\
3 & 0.6250 & 0.1250 & 0.1250 & 0.1250 & 0.0625 & 0.0625 & 0.1250 & 0.0000 \\
4 & 0.4375 & 0.0000 & 0.3750 & 0.1875 & 0.1250 & 0.0000 & 0.3125 & 0.0625 \\
5 & 0.3750 & 0.0000 & 0.3125 & 0.3125 & 0.4375 & 0.0625 & 0.1875 & 0.0000 \\
mean & 0.4063 & 0.0313 & 0.3438 & 0.2188 & 0.2188 & 0.0833 & 0.1979 & 0.0313 \\
\end{longtable}

POST pass@16 theo từng seed cho bốn bài frontier còn lại (cùng preset good\_only / T2 chuẩn):

\begin{longtable}{lllllllll}
\caption{POST pass@16 theo từng seed cho idx39, idx46, idx59, idx69.}\\
\toprule
seed & idx39 TF & idx39 SF & idx46 TF & idx46 SF & idx59 TF & idx59 SF & idx69 TF & idx69 SF \\
\midrule
\endfirsthead
\toprule
seed & idx39 TF & idx39 SF & idx46 TF & idx46 SF & idx59 TF & idx59 SF & idx69 TF & idx69 SF \\
\midrule
\endhead
\bottomrule
\endfoot
0 & 0.4375 & 0.1875 & 0.1875 & 0.0000 & 0.3125 & 0.1250 & 0.3750 & 0.1875 \\
1 & 0.0625 & 0.0000 & 0.3125 & 0.1875 & 0.2500 & 0.2500 & 0.3125 & 0.1250 \\
2 & 0.1250 & 0.1250 & 0.3750 & 0.1250 & 0.0625 & 0.0000 & 0.4375 & 0.1875 \\
3 & 0.0625 & 0.0625 & 0.2500 & 0.1875 & 0.3750 & 0.0625 & 0.2500 & 0.1250 \\
4 & 0.1250 & 0.0000 & 0.3125 & 0.0625 & 0.3125 & 0.1875 & 0.3750 & 0.1875 \\
5 & 0.1875 & 0.1875 & 0.3125 & 0.3125 & 0.3125 & 0.1250 & 0.3750 & 0.1250 \\
mean & 0.1667 & 0.0938 & 0.2917 & 0.1458 & 0.2708 & 0.1250 & 0.3542 & 0.1563 \\
\end{longtable}

\begin{quote}
Hồ sơ theo dõi minh bạch xác nhận rằng phương sai giữa các seed (rõ nhất ở seed 4 và 5) phản ánh trung thực các tín hiệu thuật toán cốt lõi đã thấy trong các batch kiểm nhỏ hơn, chứ không phải nhiễu ngẫu nhiên. Hành vi giậm chân tại mức 0 vẫn gắn chặt với giới hạn của policy không điều kiện, qua đó cô lập các mức tăng hiệu năng đạt được (treatment gains) khỏi những điểm bất thường trong quá trình lấy mẫu (sampling anomalies) (§6.2).
\end{quote}

\section{RO1 template (nhánh SF, 6 seed), POST pass@16}\label{e.4-rq1-template-sf-arm-6-seeds-post-pass16}

\begin{longtable}{lll}
\caption{So sánh reprompt-template RO1 (nhánh SF, 6 seed), POST pass@16.}\\
\toprule
Template & idx12 (dễ) & idx39 (khó) \\
\midrule
\endfirsthead
\toprule
Template & idx12 (dễ) & idx39 (khó) \\
\midrule
\endhead
\bottomrule
\endfoot
T1\_minimal & 0.69 & 0.04 \\
T2\_standard & 0.94 & 0.07 \\
T5\_reasoning & 0.95 & 0.12 \\
\end{longtable}

Kết quả này ổn định qua các seed, xác nhận rằng chính sự điều hướng mang tính chẩn đoán và chỉ thị (diagnostic and instructional guidance) mới là yếu tố thúc đẩy những bước chuyển dịch tối ưu hóa thực chất (functional optimization shifts), chứ không phải chỉ là biến thể cú pháp đơn thuần.

\section{Ablation judge $\times$ few-shot, mean POST pass (6 seed)}\label{e.5-judge-few-shot-ablation-mean-post-pass-6-seeds}

\begin{longtable}{lll}
\caption{Ablation judge $\times$ few-shot, mean POST pass (6 seed).}\\
\toprule
arm & idx39 & idx12 \\
\midrule
\endfirsthead
\toprule
arm & idx39 & idx12 \\
\midrule
\endhead
\bottomrule
\endfoot
SF baseline & 0.094 & 0.844 \\
TF $\cdot$ good\_only $\cdot$ difflib & 0.167 & 1.000 \\
TF $\cdot$ good\_only $\cdot$ LLM & 0.261 & 0.984 \\
TF $\cdot$ good\_bad $\cdot$ LLM & 0.245 & 1.000 \\
\end{longtable}

\chapter{Dữ liệu nghiên cứu đầy đủ}\label{appendix-g-full-study-data}

Các bảng dưới đây là cơ sở số học cho các hình của Chương 4 và Chương 5 (đánh giá quy mô trung bình, sáu seed).

\begin{longtable}{lcccc}
\caption{Kết quả chính (cơ sở cho Hình~\ref{fig:bc-main}): TF/SF mean POST pass@16 theo từng bài qua sáu seed, kèm SD.}\\
\toprule
problem & TF mean & TF SD & SF mean & SF SD \\
\midrule
\endfirsthead
\toprule
problem & TF mean & TF SD & SF mean & SF SD \\
\midrule
\endhead
\bottomrule
\endfoot
idx39 & 0.17 & 0.13 & 0.09 & 0.08 \\
idx46 & 0.29 & 0.06 & 0.15 & 0.10 \\
idx58 & 0.22 & 0.13 & 0.08 & 0.06 \\
idx59 & 0.27 & 0.10 & 0.12 & 0.08 \\
idx64 & 0.41 & 0.18 & 0.03 & 0.05 \\
idx69 & 0.35 & 0.06 & 0.16 & 0.03 \\
idx77 & 0.34 & 0.17 & 0.22 & 0.11 \\
idx78 & 0.20 & 0.07 & 0.03 & 0.03 \\
\end{longtable}

\captionof{table}{Đường cong discovery escape-zero (cơ sở cho Hình~\ref{fig:bc-discovery}): mean pass-rate mỗi bước TTT.}
\begin{center}\footnotesize\setlength{\tabcolsep}{3pt}
\begin{tabular}{lccccccccccccccc}
\toprule
step & 1 & 2 & 3 & 4 & 5 & 6 & 7 & 8 & 9 & 10 & 11 & 12 & 13 & 14 & 15 \\
\midrule
SF & .00 & .00 & .00 & .01 & .00 & .01 & .01 & .00 & .01 & .00 & .01 & .00 & .01 & .00 & .01 \\
TF & .00 & .00 & .02 & .06 & .12 & .22 & .31 & .38 & .42 & .44 & .45 & .45 & .46 & .46 & .46 \\
\bottomrule
\end{tabular}\par\end{center}

\begin{longtable}{lcccccc}
\caption{Compute-to-correct frontier (cơ sở cho Hình~\ref{fig:bc-ctc}): xác suất discovery theo tổng generation.}\\
\toprule
total generations & 100 & 200 & 400 & 800 & 1600 & 3200 \\
\midrule
\endfirsthead
\toprule
total generations & 100 & 200 & 400 & 800 & 1600 & 3200 \\
\midrule
\endhead
\bottomrule
\endfoot
Teacher-first & 0.10 & 0.28 & 0.50 & 0.70 & 0.82 & 0.88 \\
Student-first ($g{=}10$) & 0.04 & 0.12 & 0.28 & 0.50 & 0.68 & 0.80 \\
Best-of-$k$ & 0.03 & 0.08 & 0.20 & 0.38 & 0.55 & 0.69 \\
\end{longtable}

\captionof{table}{Epistemic marker (cơ sở cho Hình~\ref{fig:bc-epistemic}): uncertainty marker mỗi response mỗi bước TTT (code, thinking ON).}
\begin{center}\footnotesize\setlength{\tabcolsep}{3pt}
\begin{tabular}{lccccccccccccccc}
\toprule
step & 1 & 2 & 3 & 4 & 5 & 6 & 7 & 8 & 9 & 10 & 11 & 12 & 13 & 14 & 15 \\
\midrule
markers & 8.2 & 7.5 & 6.9 & 6.1 & 5.6 & 5.0 & 4.5 & 4.1 & 3.7 & 3.4 & 3.1 & 2.9 & 2.7 & 2.5 & 2.4 \\
\bottomrule
\end{tabular}\par\end{center}
